\loai{Bài toán một lượng khí thoát khỏi bình}
\begin{vidu} %[1P5-5.3B][Loai2]
	Một bình chứa khí ở nhiệt độ $27\doC$ và áp suất 4\ atm. Nếu $\dfrac{1}{2}$ khối khí thoát ra khỏi bình và nhiệt độ hạ xuống tới $12\doC$ thì khí trong bình còn lại sẽ có áp suất là bao nhiêu?
	
	\choice
	{2\ atm}
	{\True 1,9\ atm}
	{1,5\ atm}
	{2,1\ atm}
	
	\loigiai{
		\begin{itemize}
			\item Theo phương trình Mendeleep – Claperon: ${p_1}V =nR{T_1}$ (1)
			${p_2}V = 0,5nR{T_2}$ (2)
			\item Chia hai vế của phương trình (1) và (2) cho nhau $\Rightarrow \dfrac{p_1}{p_2} = \dfrac{2T_1}{T_2} \Rightarrow {p_2} = \dfrac{p_1T_2}{2T_1} = \dfrac{4.285}{2.300} = 1,9\ atm$.
		\end{itemize}
		
	}
\end{vidu}
\begin{vidu} %[1P5-5.3G][Loai2]
	Một bình chứa m = 0,3\ kg hêli. Sau một thời gian do bị hở, khí hêli thoát ra một phần. Nhiệt độ tuyệt đối của khí giảm tới 10\%, áp suất giảm 20\%. Tính số nguyên tử hêli đã thoát ra khỏi bình
	
	\choice
	{$45,2.10^{22}$ nguyên tử}
	{$5,02.10^{23}$ nguyên tử}
	{$4,52.10^{22}$ nguyên tử}
	{\True $50,2.10^{23}$ nguyên tử}
	
	\loigiai{
		\begin{itemize}
			\item Ban đầu, khí hêli có khối lượng m, thể tích V, áp suất $p_1$, nhiệt độ $T_1$: ${p_1}V = \dfrac{m}{M}R{T_1}$ (1)
			\item Sau một thời gian, khí hêli có khối lượng m', thể tích V, áp suất $p_2$, nhiệt độ $T_2$: ${p_2}V = \dfrac{m'}{M}R{T_2}$ (2)
			\item Từ (1) và (2) $\Rightarrow$ $\dfrac{p_2}{p_1} = \dfrac{m'}{m}\dfrac{T_2}{T_1}$ $\Leftrightarrow \dfrac{{{p_2} - {p_1}}}{p_1} = \dfrac{{m'{T_2} - m{T_1}}}{m{T_1}} = \dfrac{{m'\left( {{T_1} + \Delta T} \right) - m{T_1}}}{m{T_1}}\Rightarrow \dfrac{\Delta p}{p_1} = \dfrac{m' - m}{m} + \dfrac{m'}{m}\dfrac{\Delta T}{T_1}$
			\item Với $\dfrac{\Delta p}{p_1} =  - 0,2$, $\dfrac{\Delta T}{T_1} =  - 0,1$ $\Rightarrow - 0,2 = \dfrac{m'}{m} - 1 + \dfrac{m'}{m}\left( {- 0,1} \right) = \dfrac{0,9m'}{m} - 1$\\ $\Rightarrow m' = \dfrac{8}{9}m\Rightarrow \Delta m = m - m' = \dfrac{m}{9} = \dfrac{0,3}{9} = \dfrac{1}{30}\ kg = 33,333\ g$
			\item Số nguyên tử hêli đã thoát khỏi bình là $\Delta N = \dfrac{\Delta m}{M}.{N_A} = \dfrac{33,333}{4}.6,023.10^{23} = 50,2.10^{23}$
		\end{itemize}
		
	}
\end{vidu}

\loai{Bài toán đồ thị đẳng quá trình}
\begin{vidu} %[1P5-5.3K][Loai4]
	Một lượng khí hêli có khối lượng $m = 1\ g$, nhiệt độ ${t_1} = 127\doC$ và thể tích $V_1 = 4$ lít biến đổi qua hai giai đoạn: đẳng nhiệt, thể tích tăng gấp hai lần, sau đó đẳng áp, thể tích trở về giá trị ban đầu. Tìm nhiệt độ và áp suất thấp nhất trong quá trình biến đổi.
	
	\choice
	{1,05\ atm; $200\doC$}
	{\True 1,05\ atm; $-73\doC$}
	{2,1\ atm; $200\doC$}
	{2,1\ atm; $-73\doC$}
	
	\loigiai{
		\begin{align*}\begin{cases} {p_1}\\ {V_1} = 4\ \ell \\ {T_1} = 400\ K  \end{cases} \to \begin{cases} {p_2}\\ {V_2} = 2{V_1} = 8\ \ell \\ {T_2} = {T_1} = 400\ K  \end{cases}\to \begin{cases} {p_3} = {p_2}\\ {V_3} = {V_4} = 4\ \ell \\ {T_3}  \end{cases}\end{align*}
		immini[thm]{\begin{itemize}
				\item Ta có: ${p_1} = \dfrac{m}{M}\dfrac{RT_1}{V_1} = \dfrac{1}{4}\dfrac{8,31.400}{4.10^{- 3}} = 2,1.10^5\ Pa = 2,1\ atm$
				\item Từ đồ thị $\Rightarrow$ $p_{\min} = {p_2}$ 
				\item Xét quá trình đẳng nhiệt (1) – (2), ta có: ${p_1}{V_1} = {p_2}{V_2} \Rightarrow {p_2} = \dfrac{p_1V_1}{V_2} = \dfrac{2,1.4}{8} = 1,05\ atm$
				\item Từ đồ thị $\Rightarrow T_{\min} = {T_3}$ 
				\item Xét quá trình đẳng áp (2) – (3), ta có: $\dfrac{V_3}{V_2} = \dfrac{T_3}{T_2} \Rightarrow {T_3} = \dfrac{V_3}{V_2}{T_2} = \dfrac{4}{8}.400 = 200\ K$ hay ${T_3} =  - 73\doC$.
		\end{itemize}}{\begin{tikzpicture}[scale=1,thick,>=stealth']
				\tkzDefPoints{0/0/o, 4/3/a, 4/1.5/b, 2/1.5/c}
				\draw[->] (0,0)--(5,0)node[below]{$T$}; 
				\draw[->] (0,0)--(0,3.5)node[left]{$p$};
				\draw[dashed](0,1.5)--(c)--(2,0)(0,3)--(a)(b)--(4,0)(o)--(a);
				\tkzDrawSegments[blue,very thick,arrow data={0.5}{stealth}](a,b b,c)
				\node at (a)[above right]{\small(1)};
				\node at (b)[right]{\small(2)};
				\node at (c)[above left]{\small (3)};
				\node at (0,3)[left]{\small $p_{1}$};
				\node at (0,1.5)[left]{\small $p_{2}$};
				\node at (2,0)[below]{\small $T_{3}$};
				\node at (4,0)[below]{\small $T_{1}$};
		\end{tikzpicture}}
	}
\end{vidu}

\loai{Bài toán hai bình khí}
\begin{vidu} %[1P5-5.3K][Loai5]
	Hai bình có thể tích $V_1 = 100\ cm^3$, $V_2  = 200\ cm^3$ được nối bằng một ống nhỏ cách nhiệt. Ban đầu hệ có nhiệt độ $t = 27\doC$ và chứa oxygen ở áp suất $p = 760\ mmHg$. Sau đó bình $V_1$ được giảm nhiệt độ xuống $0\doC$ còn bình $V_2$ tăng nhiệt độ lên đến $100\doC$. Tính áp suất khí trong các bình khi đó.
	
	\choice
	{${p_1}' = {p_2}' = 598\ mmHg$}
	{${p_1}' = {p_2}' = 399\ mmHg$}
	{${p_1}' = {p_2}' = 1198\ mmHg$}
	{\True ${p_1}' = {p_2}' = 842\ mmHg$}
	
	\loigiai{
		\begin{itemize}
			\item Ban đầu: 
			\begin{itemize}\item Bình I có thể tích ${V_1} = V$, áp suất p, nhiệt độ T 
				\item Bình II có thể tích ${V_2} = 2V$, áp suất p, nhiệt độ T.
				\item Tổng số mol khí trong hai bình là $n = \dfrac{p.3V}{RT}$
			\end{itemize}
			\item Sau đó: 
			\begin{itemize}\item Bình I có thể tích $V_1$, áp suất p', nhiệt độ $T_1$ nên số mol khí của bình I là ${n_1} = \dfrac{p'V}{RT_1}$
				\item Bình II có thể tích $V_2$, áp suất p', nhiệt độ $T_2$ nên số mol khí của bình II là ${n_2} = \dfrac{p'2V}{RT_2}$
			\end{itemize}
			\item Mặt khác $n = {n_1} + {n_2} \Leftrightarrow \dfrac{p.3V}{RT} = \dfrac{p'V}{RT_1} + \dfrac{p'2V}{RT_2}$\\$\Leftrightarrow \dfrac{3p}{T} = p'\left( {\dfrac{1}{T_1} + \dfrac{2}{T_2}} \right) \Rightarrow p' = \dfrac{3p}{{T\left( {\dfrac{1}{T_1} + \dfrac{2}{T_2}} \right)}} = \dfrac{3.760}{{300\left( {\dfrac{1}{273} + \dfrac{2}{373}} \right)}} = 842\ mmHg$.
		\end{itemize}
	}
\end{vidu}

\loai{Bài toán tìm $T_{max}$ từ đồ thị}
\begin{vidu} %[1P5-5.3T][Loai7] 
	immini[thm]{Có 20\ g khí hêli chứa trong xilanh đậy kín bởi pít-tông biến đổi chậm từ (1) - (2) theo đồ thị như hình vẽ. Cho $V_1=30$ lít; $p_1=5$ atm; $V_2=10$ lít; $p_2=15$ atm. Tìm nhiệt độ cao nhất mà khí đạt được trong quá trình trên}{\begin{tikzpicture}[scale=0.8,thick,>=stealth']
			\tkzDefPoints{0/0/o, 3/1/a, 1/3/b}
			\draw[->] (0,0)--(4,0)node[below]{$V$}; 
			\draw[->] (0,0)--(0,4)node[left]{$p$};
			\draw[dashed](0,3)--(b)--(1,0)(0,1)--(a)--(3,0);
			\tkzDrawSegments[blue,very thick,arrow data={0.5}{stealth}](a,b)
			\node at (a)[right]{\small(1)};
			\node at (b)[above]{\small(2)};
			\node at (0,1)[left]{\small $p_{2}$};
			\node at (0,3)[left]{\small $p_{1}$};
			\node at (1,0)[below]{\small $V_{2}$};
			\node at (3,0)[below]{\small $V_{1}$};
	\end{tikzpicture}}
	
	\loigiai{
		\begin{itemize}
			\item Quá trình (1) - (2): $p=aV+b$.
			\item Thay các giá trị $(p_1;\ V_1)$, $(p_2;\ V_2)$ vào (1) ta được:
			\begin{align*}
				\left\{\begin{aligned}& 5=30a+b\quad\quad(1) \\
					& 10=10a+b\quad\quad(2)
				\end{aligned}\right.
			\end{align*}
			\item Từ (1) và (2) suy ra $a=-\dfrac{1}{2}$; $b=20$ $\Rightarrow p=-\dfrac{V}{2}+20\Rightarrow pV=-\dfrac{V^2}{2}+20V\quad\quad(3)$.
			\item Mặt khác: $pV=\dfrac{m}{\mu}RT=\dfrac{20}{4}RT=5RT\quad\quad(4)$.
			\item Từ (4) suy ra: $T=-\dfrac{V^2}{10R}+\dfrac{4V}{R}\quad\quad(5)$.
			\item Xét hàm số $T=f(V)$. Ta có $T=T_{max}$ khi $V=-\dfrac{\dfrac{R}{4}}{2\left(-\dfrac{1}{10R}\right)}=20$ lít, lúc đó
			\begin{align*}
				T_{max}=-\dfrac{20^2}{10.0,082}+\dfrac{4.20}{0,082}=487,8\ K.
			\end{align*}
			\item Vậy nhiệt độ cao nhất mà khí đạt được trong quá trình biến đổi là 487,8\ K.
		\end{itemize}
	}
\end{vidu}

\loai{Bài toán khác}
\begin{vidu}%[1P5-5.3K][Loai7]
	Có $1\ g$ oxygen ở áp suất $3\ atm$ sau khi hơ nóng đẳng áp nó chiếm thể tích $1\text{ lít}$. Tìm nhiệt độ của oxygen sau khi hơ nóng. Coi khí oxygen là khí lí tưởng.
	\loigiai{
		\begin{itemize}
			\item Áp dụng định luật Gay Luy-xác cho quá trình đẳng áp:
			\begin{align*}
				\dfrac{V_1}{T_1}=\dfrac{V_2}{T_2}\Rightarrow T_2=\dfrac{V_2}{V_1}T_1\tag{1}
			\end{align*}
			\item Áp dụng phương trình trạng thái cho trạng thái 1:
			\begin{align*}
				p_1V_1=\dfrac{m}{\mu}\tag{2}
			\end{align*}
			\item Từ $(1)$ và $(2)$ rút ra:
			\begin{align*}
				T_2=\dfrac{\mu p_1V_2}{mR}
			\end{align*}
			\item Thay số $\mu=32\ g/mol=32.10^{-3}\ kg/mol$, $p_1=3\ atm=3.9,81.10^4\ N/m^2;\ V_2=1\text{ lít}=10^{-3}\ m^3$, ta tìm được: $T_2=1133\ K$
		\end{itemize}
	}
\end{vidu}
\begin{vidu} %[1P5-5.3K][Loai6]
	Trong một ống dẫn khí tiết diện đều $S = 5\ cm^2$ có khí $CO_2$ chảy qua ở nhiệt độ $35\doC$ và áp suất $3.10^5\ N/m^2$. Tính vận tốc của dòng khí biết trong thời gian 10 phút có m = 3\ kg khí $CO_2$ qua tiết diện ống.
	
	\choice
	{1,588\ m/s}
	{1,399\ m/s}
	{1,858\ m/s}
	{\True 1,939\ m/s}
	
	\loigiai{
		\begin{itemize}
			\item Thể tích khí qua ống trong thời gian 10 phút là: $V = vSt$ 
			\item Áp dụng phương trình Mendeleep – Claperon: \begin{align*}pV = \dfrac{m}{M}RT \Rightarrow pvSt = \dfrac{mRT}{M} \Rightarrow v = \dfrac{mRT}{MpSt} = \dfrac{3000.8,31.308}{44.3.10^5.5.10^{- 4}.10.60} = 1,939\ m/s\end{align*}
		\end{itemize}
	}
\end{vidu}



